\subsubsection{Alternating Operators and \texorpdfstring{$p$}{p} Layers}\label{sec:alternating-operators-and-p-layers}

Until now, we have seen the two main operators of QAOA independently:
\begin{itemize}
    \item The \textbf{Cost Operator}:
    \begin{equation*}
        U_C(\gamma) = e^{-i \gamma H_C}
    \end{equation*}
    \item The \textbf{Mixer Operator}:
    \begin{equation*}
        U_M(\beta) = e^{-i \beta H_M}
    \end{equation*}
\end{itemize}
The natural question now is: \emph{\textbf{how are these operators actually used together to build QAOA?}} And the answer is very simple: they are \textbf{applied alternately, layer after layer}. This alternating structure is exactly what gives the algorithm its name: \emph{Quantum Approximate Optimization Algorithm}.

\highspace
\textcolor{Green3}{\faIcon{tools} \textbf{One QAOA layer.}} Let's recall the QAOA circuit from \autoref{fig:qaoa-ansatz} (\autopageref{fig:qaoa-ansatz}). One \textbf{QAOA layer} is simply the application of the cost operator followed by the mixer operator, as shown in \autoref{fig:qaoa-ansatz}:
\begin{equation*}
    U_C\left(\gamma_i\right) = e^{-i \gamma_i H_C} \quad \text{and} \quad U_M\left(\beta_i\right) = e^{-i \beta_i H_M}
\end{equation*}
Mathematically, a single QAOA layer can be expressed as:
\begin{equation}\label{eq:qaoa-layer}
    U(\gamma_i, \beta_i) = U_M(\beta_i) U_C(\gamma_i) = e^{-i \beta_i H_M} \cdot e^{-i \gamma_i H_C}
\end{equation}
The ansatz for a single layer of QAOA with two qubits is shown in the following circuit diagram:
\begin{figure}[!htp]
    \centering
    \begin{quantikz}
        \lstick{\ket{0}} & \gate{H} & \gate{U_C(\gamma_1)}\gategroup[2,steps=1,style={inner sep=2pt}]{$U_C(\gamma_i)$}  & \gate{U_M(\beta_1)}\gategroup[2,steps=1,style={inner sep=2pt}]{$U_M(\beta_i)$}    & \meter{} & \rstick{\dots} \\
        \lstick{\ket{0}} & \gate{H} & \gate{U_C(\gamma_2)}                                                              & \gate{U_M(\beta_2)}                                                               & \meter{} & \rstick{\dots}
    \end{quantikz}
    \caption{QAOA ansatz with one layer ($p=1$) for two qubits ($n=2$).}
\end{figure}

\noindent
\textcolor{Green3}{\faIcon{question-circle} \textbf{Why this order?}} Initially, the algorithm prepares $\ket{+}^{\otimes n}$, which is an equal superposition of all possible solutions.\footnote{%
    Actually, the algorithm prepares the state $\ket{0}^{\otimes n}$ and then applies a Hadamard gate to each qubit to create the superposition.%
} At this point, every solution has the same probability, we have maximum exploration of the solution space, and no preference for good or bad solutions. Suppose we applied the mixer first; the state $\ket{+}^{\otimes n}$ is actually an eigenstate of the mixer Hamiltonian $H_M$, and thus:
\begin{equation*}
    U_M\left(\beta\right) \ket{+}^{\otimes n} = e^{-i \psi} \ket{+}^{\otimes n}
\end{equation*}
Where $e^{-i \psi}$ is only a \textbf{global phase}, which has \textbf{no physical effect}. So the first mixer would essentially do \textbf{nothing useful}.

\highspace
Instead, we first apply the cost operator, which \textbf{modifies the amplitudes of the states based on their cost}. This is the first step in guiding the algorithm towards better solutions. Then, we apply the mixer operator, which \textbf{mixes the amplitudes of the states}, allowing for exploration of new solutions while still favoring those with lower costs.

\highspace
\textcolor{Green3}{\faIcon{tools} \textbf{The complete QAOA layer.}} We can \textbf{extend} the single layer \autoref{eq:qaoa-layer} to $p$ layers, where each layer has its own parameters $\gamma_i$ and $\beta_i$. The complete QAOA circuit with $p$ layers can be expressed as:
\begin{equation}
    \ket{\psi(\vec{\gamma}, \vec{\beta})} = U_M\left(\beta_p\right) U_C\left(\gamma_p\right) \cdots U_M\left(\beta_1\right) U_C\left(\gamma_1\right) \ket{+}^{\otimes n}
\end{equation}
Where:
\begin{itemize}
    \item $\vec{\gamma} = \left(\gamma_1, \gamma_2, \ldots, \gamma_p\right)$ are the \textbf{parameters} for the \textbf{cost operators}.
    \item $\vec{\beta} =  \left(\beta_1, \beta_2, \ldots, \beta_p\right)$ are the \textbf{parameters} for the \textbf{mixer operators}.
    \item Parameter $p$ is the \textbf{depth} of the QAOA circuit, which controls the number of alternating layers of cost and mixer operators. It is defined as the number of Cost-Mixer pairs in the circuit.
\end{itemize}
